\documentclass{abntexto-ufc}

\makeatletter
\renewcommand{\ufcObjectLegendHook}{%
  \typeout{UFC-OBJECT-TITLE-WIDTH=\the\hsize}%
  \typeout{UFC-OBJECT-CONTENT-WIDTH=\the\savedplacewidth}%
  \typeout{UFC-OBJECT-TITLE-FONTSIZE=\f@size}%
}
\newcommand{\ufcSourceProbe}{%
  \typeout{UFC-OBJECT-SOURCE-WIDTH=\the\hsize}%
  \typeout{UFC-OBJECT-SOURCE-FONTSIZE=\f@size}%
}
\newcommand{\ufcNoteProbe}{%
  \typeout{UFC-OBJECT-NOTE-WIDTH=\the\hsize}%
  \typeout{UFC-OBJECT-NOTE-FONTSIZE=\f@size}%
}
\makeatother

\begin{document}

\legend{figure}{Título deliberadamente longo para validar que a legenda respeita a largura física da ilustração}
\ufcSource{\ufcSourceProbe Elaboração própria com texto deliberadamente longo para validar a largura da fonte.}
\ufcNote{\ufcNoteProbe Validação geométrica de nota deliberadamente longa em um objeto estreito.}
\begin{ufcobject}[here]
  \rule{6cm}{2cm}
\end{ufcobject}

\end{document}
